\documentclass[10pt,a4paper]{amsart}
\usepackage[margin=19mm]{geometry}
\usepackage{amsmath,amssymb,mathtools}
\usepackage[scr=rsfs,cal=euler]{mathalfa}
\usepackage{xfrac}
\usepackage[unicode,hidelinks,bookmarks=false]{hyperref}
\newcommand{\ie}{i.e.\ }
\newcommand{\cf}{cf.\ }
\newcommand{\resp}{respectively\ }
\newcommand{\Subsection}{\subsection}
\providecommand{\nobreakdash}{\nobreak\hbox{-}}
\setlength{\emergencystretch}{2em}
\multlinegap0pt
\usepackage{bm}
\usepackage{bbm}
\usepackage{dsfont}
\usepackage{xspace}
\usepackage{cancel}
\usepackage{mathtools}
\usepackage{amsthm}
\makeatletter
\@ifundefined{definition}{\newtheorem{definition}{Definition}}{}
\@ifundefined{lemma}{\newtheorem{lemma}{Lemma}}{}
\makeatother
\title[Inhomogeneous Khintchine-Groshev theorem without monotonicit]{Inhomogeneous Khintchine-Groshev theorem without monotonicity}
\author{Seongmin Kim}
\date{Source: arXiv:2411.07932v2 (CC BY 4.0)}
\begin{document}
\maketitle

\noindent\emph{Source and licence.} Inhomogeneous Khintchine-Groshev theorem without monotonicity. Version arXiv:2411.07932v2 (CC BY 4.0), distributed under the Creative Commons Attribution 4.0 International licence, as recorded in the arXiv licence metadata for this exact version and shown on the abstract page https://arxiv.org/abs/2411.07932v2. Full rights notice: licenses/arxiv-2411.07932-CC-BY-4.0.txt.\par\smallskip

\noindent\textit{Editorial scope.} Counted unit: the Lemma whose source label is \texttt{mainlemma} in the article (source0.tex lines 215--218, source label \texttt{mainlemma}), delivered together with the article's own complete original proof of it (source0.tex lines 219--247, one \texttt{proof} environment of 29 lines). No mathematics is written, repaired or re-derived: statement and proof are the article's own text line for line. Required context retained uncounted: a,b (lines 83--85); def (lines 112--119). Every definition, display and label of those lines is retained unchanged. Omitted from this package and not redistributed: 1--82, 86--111, 120--214, 248--568. Nothing inside a retained excerpt is omitted or altered: every retained body is delivered exactly as the article's lines are written, so no omission receipt of any kind accompanies this package. Citations: the retained excerpts carry no citation command at all, so no bibliography entry of the article is transcribed and no reference is redistributed. Reference adaptation: none was needed and none was made; every reference inside the delivered unit resolves inside it, so no unresolved reference or citation is produced. Printed slips kept literal: no printed word, symbol, hypothesis or number of the counted statement or of its proof was altered. Layout and class: the article's own class is \texttt{amsart} with packages including amsmath, amssymb, amsfonts, mathrsfs, url, hyperref, dsfont; the wrapper is the lane's pinned \texttt{amsart} editorial wrapper at 10pt on A4, which restates the theorem-like environments the retained text needs and adds the article's own macro definitions that the retained text needs (none), with extra wrapper packages (bm, bbm, dsfont, xspace, cancel, mathtools). The delivered numbering is the unit's own. Assets: no retained span contains a figure environment, a graphics-inclusion command, a PDF-page inclusion, a table or a TikZ picture; no image file is copied into this package, the package declares no asset dependency and no image file is redistributed with it. Licence: version arXiv:2411.07932v2 of this article is distributed under the Creative Commons Attribution 4.0 International licence, as recorded in the arXiv licence metadata for this exact version and shown on the abstract page https://arxiv.org/abs/2411.07932v2. A full rights notice is delivered at licenses/arxiv-2411.07932-CC-BY-4.0.txt. Characters: every retained excerpt is pure ASCII, so the delivered engine, XeLaTeX with its default Latin Modern text font, reports no missing character.
\begin{equation}\label{a,b}
    |b_d y - a_d| < \frac{1}{|d|}, \quad 1\leq b_d \leq |d|, \quad \gcd(a_d, b_d) = 1.
\end{equation}

\begin{definition}\label{def}
    Let $n,m \in \mathbb{N}$, and let $(\frac{\mathbf{a}_\mathbf{q}}{b_\mathbf{q}})_{\mathbf{q} \in \mathbb{Z}^n \setminus \{\mathbf{0}\}}$ be a sequence of rational vectors in $\mathbb{Q}^m$ such that $\mathbf{a}_\mathbf{q} \in \mathbb{Z}^m,\ b_\mathbf{q} \in \mathbb{N}$ and $\gcd(\mathbf{a}_\mathbf{q}, b_\mathbf{q}) = 1.$ Also, let $(\mathbf{y}_\mathbf{q})_{\mathbf{q} \in \mathbb{Z}^n \setminus \{\mathbf{0}\}}$ be a sequence of vectors in $\mathbb{R}^m$. We call $\mathbf{y}_\mathbf{q}$ the \emph{moving targets}, and when $\mathbf{y}_\mathbf{q} = \mathbf{y}$ is fixed for every $\mathbf{q}$, then we call $\mathbf{y}$ the \emph{inhomogeneous parameter}. For each $\mathbf{q} \in \mathbb{Z}^n \setminus \{\mathbf{0}\}$ and $\delta \in \mathbb{R}_{\geq 0},$ we define the following sets:
\begin{align*}
    A_{n,m}(\mathbf{q},\delta) &= \{\mathbf{x} \in [0,1)^{nm} : \exists \mathbf{p} \in \mathbb{Z}^m \text{ s.t. } |\mathbf{qx} - \mathbf{p} - \mathbf{y}_\mathbf{q}| < \delta\}, \\
    \tilde{A}_{n,m}(\mathbf{q},\delta) &= \{\mathbf{x} \in [0,1)^{nm} : \exists \mathbf{p} \in \mathbb{Z}^m \text{ s.t. } |\mathbf{qx} - \mathbf{p} - \mathbf{y}_\mathbf{q}| <\delta, \ \gcd(\mathbf{q}, b_\mathbf{q}\mathbf{p} + \mathbf{a}_\mathbf{q}) = 1\}.
\end{align*}
    When a \emph{multivariate function} $\Psi : \mathbb{Z}^n \setminus \{\mathbf{0}\} \to \mathbb{R}_{\geq 0}$ is given, we denote $A_\mathbf{q} := A_{n,m}(\mathbf{q}, \Psi(\mathbf{q})),\ \tilde{A}_\mathbf{q} := \tilde{A}_{n,m}(\mathbf{q}, \Psi(\mathbf{q}))$. Similarly, when a \emph{univariate function} $\psi : \mathbb{N} \to \mathbb{R}_{\geq 0}$ is given, we denote $A_\mathbf{q} := A_{n,m}(\mathbf{q}, \psi(|\mathbf{q}|)),\ \tilde{A}_\mathbf{q} := \tilde{A}_{n,m}(\mathbf{q}, \psi(|\mathbf{q}|))$.
\end{definition}

\begin{lemma}\label{mainlemma}
    Let $\psi: \mathbb{N} \to [0,\frac{1}{2})$ satisfy $\psi(q) \leq q^{-1}$ for all $q \in \mathbb{N},$ and let $y \in \mathbb{R}$ denote an inhomogeneous parameter. Define $(a_d,b_d)_{d \in \mathbb{Z} \setminus \{0\}}$ as in (\ref{a,b}), and  $\tilde{A}$ as in Definition \ref{def}. Let integers $q,r,d,e$ satisfy $1\leq r < q$ and $1\leq |e| < |d|$. If $\gcd(d,e) \geq 3$ and $r \geq 2d^2,$ then
    $$|\tilde A_{1,1} (d, \psi(q)) \cap \tilde A_{1,1}(e, \psi(r))| = 0.$$
\end{lemma}

\begin{proof}
    We are going to show that $\tilde A_{1,1} (d, \psi(q)) \cap A_{1,1}(e, \psi(r)) = \emptyset$ which is a stronger result. The definitions of the sets $\tilde A_{1,1}(d, \psi(q))$ and $A_{1,1}(e, \psi(r))$ give
    \begin{equation}\label{union}
        \begin{aligned}
            \tilde A_{1,1}(d, \psi(q)) &= \bigcup_{\substack{p \in \mathbb{Z}\\ \gcd(d, b_d p + a_d)=1}} \left(\frac{p+y}{d} -\frac{\psi(q)}{|d|}, \frac{p+y}{d} + \frac{\psi(q)}{|d|} \right) \cap [0,1),\\
            A_{1,1}(e, \psi(r)) &= \bigcup_{s \in \mathbb{Z}} \left(\frac{s+y}{e} - \frac{\psi(r)}{|e|}, \frac{s+y}{e} + \frac{\psi(r)}{|e|} \right) \cap [0,1).
        \end{aligned}
    \end{equation}
    
    The distance between the centers of intervals in the unions (\ref{union}) is
    \begin{align}\label{distance}
        \left| \frac{p+y}{d} - \frac{s+y}{e} \right| &= \left| \frac{p+ \frac{a_d}{b_d} +y -\frac{a_d}{b_d}}{d} - \frac{s+ \frac{a_d}{b_d} + y - \frac{a_d}{b_d}}{e} \right| \nonumber \\
        &= \left| \frac{1}{b_d} \left(\frac{b_d p + a_d}{d} - \frac{b_d s + a_d}{e} \right) + \left(y - \frac{a_d}{b_d} \right) \left(\frac{1}{d} - \frac{1}{e} \right)\right|.
    \end{align}
    Since $\gcd(d, b_d p+a_d) = 1$ and $|e| < |d|$, $\frac{b_d p+a_d}{d}$ is a reduced fraction with a denominator larger than that of $\frac{b_d s+a_d}{e}.$ Thus the two fractions $\frac{b_d p+a_d}{d}$ and $\frac{b_d s+a_d}{e}$ are distinct, so the first term of (\ref{distance}) satisfies
    $$\left| \frac{1}{b_d} \left(\frac{b_d p + a_d}{d} - \frac{b_d s + a_d}{e} \right) \right| \geq \frac{1}{b_d \text{lcm}(d,e)} \geq \frac{3}{b_d |de|},$$
    where the last inequality comes from $\gcd(d,e) \geq 3$. By the definition of $(a_d,b_d)$, we have the inequality $|b_d y - a_d| < \frac{1}{|d|}$. Since $|e|<|d|$, we also have $|e-d|<2|d|$, thus the second term of (\ref{distance}) satisfies
    $$\left|\left(y - \frac{a_d}{b_d} \right) \left(\frac{1}{d} - \frac{1}{e} \right)\right| < \frac{1}{b_d|d|} \frac{|e-d|}{|de|} < \frac{2}{b_d |de|}.$$
    Therefore, from (\ref{distance}) we have
    $$\left| \frac{p+y}{d} - \frac{s+y}{e} \right| > \frac{1}{b_d |de|}.$$

    We now compare the distance between centers with the sum of the radii of intervals in the unions (\ref{union}), given by $\frac{\psi(q)}{|d|} + \frac{\psi(r)}{|e|}$. Since $\psi(q) \leq q^{-1}$, $2d^2 \leq r<q,$ and $|e|<|d|$, it follows that
    $$\frac{\psi(q)}{|d|} + \frac{\psi(r)}{|e|} \leq \frac{1}{q|d|} + \frac{1}{r|e|} < \frac{2}{r|e|} \leq \frac{1}{d^2|e|}.$$
    By the definition of $b_d$, we have $b_d \leq |d|$. Therefore
    $$\frac{\psi(q)}{|d|} + \frac{\psi(r)}{|e|} < \frac{1}{b_d |de|} < \left| \frac{p+y}{d} - \frac{s+y}{e} \right|$$
    for any $p,s \in \mathbb{Z}$ such that $\gcd(d, b_dp+a_d) = 1$. This implies that
    $$\left(\frac{p+y}{d} -\frac{\psi(q)}{|d|}, \frac{p+y}{d} + \frac{\psi(q)}{|d|} \right) \cap \left(\frac{s+y}{e} - \frac{\psi(r)}{|e|}, \frac{s+y}{e} + \frac{\psi(r)}{|e|} \right) = \emptyset,$$
    which allows us to conclude that $\tilde A_{1,1} (d, \psi(q)) \cap A_{1,1}(e, \psi(r)) = \emptyset$.
\end{proof}

\end{document}
